\documentclass[11pt,a4paper]{article}

\usepackage{lmodern}
\usepackage[T1]{fontenc}
\usepackage[utf8]{inputenc}
\usepackage[margin=2.6cm]{geometry}
\usepackage{amsmath}
\usepackage{amssymb}
\usepackage{booktabs}
\usepackage{microtype}
\usepackage{xcolor}
\usepackage[round]{natbib}
\usepackage[hidelinks]{hyperref}

\newcommand{\code}[1]{\texttt{#1}}
\DeclareMathOperator{\sgn}{sign}
\DeclareMathOperator{\ewm}{ewm}
\DeclareMathOperator*{\argmax}{arg\,max}

\title{The Ten-Line CTA: What Five Structural Refinements Actually Buy}

\author{Thomas Schmelzer\thanks{\url{https://github.com/tschm/cs}}}

\date{\today}

\begin{document}

\maketitle

\begin{abstract}
This repository contains a trend-following strategy written in ten lines of
code and four successive refinements of it, each a marimo notebook whose
sliders expose its free parameters. This paper measures what the sequence
buys. The headline Sharpe ratio rises from $0.56$ to $1.47$ over the five
systems on the shipped panel of $49$ futures and $53$ years of daily data, but
it does not rise monotonically: the third step, which introduces the Huber
clip and the analytically scaled oscillator, leaves the Sharpe ratio unchanged
to three decimals ($0.8794 \to 0.8776$) and would be read as a null result by
anyone looking only at that number. It is not a null result. It cuts annual
one-way turnover by a factor of three, and net of a one-way cost of $5$ basis
points it is worth $+0.14$ Sharpe over the step it replaces. Turnover is in
fact where the sequence is least flattering to itself: it rises $35$-fold from
the first system to the last, so the ranking of the five is cost-dependent,
and beyond roughly $12$ basis points the ten-line original beats the
correlation-optimised finale. Against that, we find the structural steps
robustly dominate parameter search. Sweeping the two moving-average lengths
over all $852$ admissible pairs, the \emph{best} pair for any system is worse
than the \emph{default} parameters of the system two steps later, and the
Spearman rank correlation between the parameter surfaces of any two systems is
between $0.93$ and $0.98$ --- the structural change moves the level, not the
shape. Optuna searches over the same spaces add $+0.06$ to $+0.23$ Sharpe,
against $+0.32$, $+0.19$ and $+0.39$ for the three structural steps that pay.
\end{abstract}

\section{Introduction}

The premise of this repository is a challenge: write a hedge-fund strategy in
ten lines of code. The answer given is a sign-of-a-crossover trend follower, a
construction old enough that its literature is about how long it has worked
\citep{hurst2017} rather than whether it works \citep{moskowitz2012}. It is
also, in the form first written, indefensible in several specific ways that
the notebook itself lists: no notion of risk, a bet size independent of the
asset, and a signal whose only values are $\pm 1$.

What follows in the repository is not a defence of the ten lines but four
rewrites, each addressing one flaw and each still small enough to read in one
screen. The sequence is deliberately pedagogical: volatility scaling, then a
price filter and a scale-free oscillator, then cross-sectional risk scaling,
then a full covariance-aware allocation. That last step is convex
optimisation in the sense of \citet{schmelzer2013}: the position is the
solution of a small linear system involving an estimated correlation matrix,
and the parameter-estimation care that paper argues for is what makes it
usable.

This paper asks the question the notebooks leave open. Each step
\emph{improves the Sharpe ratio} --- that is what the printed number at the
bottom of each notebook shows, and it is what the sequence is arranged to
demonstrate. But a Sharpe ratio computed gross of costs on a single panel is
a weak instrument, and three questions are more interesting than the ranking
it produces:

\begin{itemize}
  \item What does each step cost in turnover, and does the ranking survive
        trading costs?
  \item How much of the improvement could have been had by tuning the
        parameters of the previous step instead --- the
        ``parameter-hacking'' the notebooks explicitly disclaim?
  \item Which of the five steps are doing something the headline number does
        not show?
\end{itemize}

Section~\ref{sec:method} states the five systems as formulas.
Section~\ref{sec:impl} records the implementation choices that are not implied
by those formulas and matter for reproducing the numbers.
Section~\ref{sec:results} answers the three questions. The short version of
the answers: costs reorder the five and can invert the sequence entirely;
parameter search is worth roughly half a structural step and never more; and
the step that looks worthless is the one that makes the rest affordable.

\section{The five systems}
\label{sec:method}

Let $p_{t,i}$ be the price of asset $i$ at time $t$, and let
$\ewm_m$ denote an exponentially weighted mean with centre of mass
$m$ (so decay $1 - 1/(m+1)$ per step), with $\ewm^{\sigma}_m$ the
matching standard deviation. Every system maps prices to a \emph{cash
position} $x_{t,i}$, which is then run through the portfolio accounting of
\code{jquantstats} against an AUM of $10^8$. Long-only is never imposed and
no leverage limit is applied; the scaling constant in each system is chosen so
that the resulting portfolio volatility is of a comparable order.

\paragraph{CTA 1.0 --- the ten lines.}
\begin{equation}
  x_{t,i} = 5 \cdot 10^{6} \cdot
  \sgn\!\left( \ewm_{32}(p_{\cdot,i})_t - \ewm_{96}(p_{\cdot,i})_t \right).
  \label{eq:cta1}
\end{equation}
Each asset carries $\pm 5\%$ of AUM in notional, whatever it is. The position
is piecewise constant and changes only when the crossover flips sign.

\paragraph{CTA 2.0 --- volatility scaling.}
\begin{equation}
  x_{t,i} = 10^{5} \cdot
  \frac{\sgn\!\left( \ewm_{32}(p_{\cdot,i})_t - \ewm_{96}(p_{\cdot,i})_t \right)}
       {\ewm^{\sigma}_{32}\!\left( \Delta p / p \right)_{t,i}} .
  \label{eq:cta2}
\end{equation}
The bet is now inversely proportional to the asset's estimated volatility, so
a bond future and an equity future are sized to contribute comparably. This is
the step that turns a signal into a portfolio. Note what it does to trading:
the numerator still only flips, but the denominator moves every day, so the
position moves every day.

\paragraph{CTA 3.0 --- filter the price, scale the signal.}
Equation~\eqref{eq:cta2} feeds a raw price into a function that is only
scale-free because $\sgn$ discards magnitude. To use anything richer than a
sign, the price has to be standardised first. Define clipped,
volatility-adjusted log returns
\begin{equation}
  \tilde r_{t,i} =
  \mathrm{clip}_{c}\!\left(
    \frac{\Delta \log p_{t,i}}{\ewm^{\sigma}_{\nu-1}(\Delta \log p_{\cdot,i})_t}
  \right),
  \qquad
  \mathrm{clip}_{c}(z) = \max(-c, \min(c, z)),
  \label{eq:voladj}
\end{equation}
with $c = 4.2$ by default, and integrate them into an adjusted price
$\tilde p_{t,i} = \sum_{s \le t} \tilde r_{s,i}$. The clip is the derivative
of the Huber loss \citep{huber1964} applied to the standardised score, which
is what makes ``we apply Huber functions'' in the notebook an accurate
description of an operation that looks like plain winsorising.

On the standardised price, the oscillator
\begin{equation}
  \mathrm{osc}(\tilde p; f, s)_t =
  \frac{\ewm_{f-1}(\tilde p)_t - \ewm_{s-1}(\tilde p)_t}{\varsigma(f,s)},
  \qquad
  \varsigma(f,s)^2 = \frac{1}{1-\phi^2} - \frac{2}{1-\phi\gamma} + \frac{1}{1-\gamma^2},
  \label{eq:osc}
\end{equation}
with $\phi = 1 - 1/f$ and $\gamma = 1 - 1/s$, divides the crossover by its
theoretical standard deviation under a unit-variance random walk. The point of
$\varsigma$ is that the oscillator has unit scale for \emph{any} pair
$(f,s)$, so no separate volatility lookback is needed to compare a fast
crossover with a slow one. The system is then
\begin{equation}
  x_{t,i} = 10^{5} \cdot
  \frac{\tanh\!\left( \mathrm{osc}(\tilde p_{\cdot,i}; f, s)_t \right)}
       {\ewm^{\sigma}_{s}\!\left( \Delta p / p \right)_{t,i}} .
  \label{eq:cta3}
\end{equation}
The $\tanh$ bounds the signal, so a large oscillator reading no longer
translates into an unbounded position.

\paragraph{CTA 4.0 --- cross-sectional risk scaling.}
Write $\mu_{t} \in \mathbb{R}^{N}$ for the vector of signals
$\mu_{t,i} = \tanh(\mathrm{osc}(\tilde p_{\cdot,i}; f, s)_t)$. The first three
systems treat the assets one at a time; this one normalises across them,
\begin{equation}
  x_{t,i} = 5 \cdot 10^{5} \cdot
  \frac{1}{\ewm^{\sigma}_{\nu}(\Delta p / p)_{t,i}} \cdot
  \frac{\mu_{t,i}}{\lVert \mu_t \rVert_2},
  \label{eq:cta4}
\end{equation}
so the total signal strength spent on any day is constant: conviction is
redistributed between assets rather than added to the book. This is the
smallest possible step from a univariate to a portfolio construction --- it
uses the cross-section but not the covariance.

\paragraph{CTA 5.0 --- correlation-aware allocation.}
The last system replaces the Euclidean norm in \eqref{eq:cta4} by one induced
by an estimated correlation matrix. Per-day correlations come from an
exponentially weighted covariance of the standardised returns
$\tilde r$ in the manner of \citet{engle2002},
\begin{equation}
  \Sigma_{t,ij} = \ewm_{\kappa}(\tilde r_i \tilde r_j)_t - \ewm_{\kappa}(\tilde r_i)_t \, \ewm_{\kappa}(\tilde r_j)_t,
  \qquad
  C_{t,ij} = \frac{\Sigma_{t,ij}}{\sqrt{\Sigma_{t,ii} \Sigma_{t,jj}}},
\end{equation}
shrunk towards the identity \citep{ledoit2004} with intensity $\lambda$,
\begin{equation}
  \tilde C_t = (1-\lambda) C_t + \lambda I .
  \label{eq:shrink}
\end{equation}
Restricting to the assets with a live price on day $t$, the position is
\begin{equation}
  x_{t} = 10^{6} \cdot D_t^{-1}
  \frac{\tilde C_t^{-1} \mu_t}{\sqrt{\mu_t^{\top} \tilde C_t^{-1} \mu_t}},
  \qquad
  D_t = \operatorname{diag}\!\left( \ewm^{\sigma}_{\nu}(\Delta p / p)_{t,\cdot} \right).
  \label{eq:cta5}
\end{equation}
The numerator is the direction that maximises signal per unit of estimated
risk; the denominator normalises that risk to one. Days on which the solve
fails, or the norm vanishes, are left flat rather than allowed to raise an
exception out of the loop.

Equation~\eqref{eq:cta5} is the only system that inverts anything, and it is
the reason the shrinkage parameter exists: the inversion is the step
\citet{michaud1989} warns about and \citet{schmelzer2013} treat as a
numerical-hygiene problem. Nothing here estimates expected returns, so the
comparison with mean--variance optimisation \citep{markowitz1952} is only
partial --- $\mu_t$ is a bounded signal, not a return forecast.

\section{Implementation}
\label{sec:impl}

\paragraph{The notebooks are the source of truth.}
Each system lives in one marimo notebook \citep{marimo} that defines a signal
function \code{f} and builds a portfolio from it. The Optuna driver
\citep{akiba2019} does not reimplement the strategies: it executes each
notebook with \code{runpy} and pulls \code{f} out of the resulting namespace,
so a parameter search cannot silently drift away from the strategy it claims
to be searching. The five headline Sharpe ratios are pinned in the test suite
to a relative tolerance of $10^{-6}$ and checked from both directions --- once
by executing the notebooks, once by calling the driver's builders --- so the
numbers in Section~\ref{sec:results} are regression-tested, not just recorded.

\paragraph{Warm-up.}
The EWM estimators are given \code{min\_samples} of $100$ (CTA 1.0) and $300$
(all later systems). With $13{,}909$ daily rows the warm-up is immaterial to
the reported statistics, but it is the reason a null-safe fill is applied
before the portfolio is built: the first rows of every signal are undefined,
and they are set to zero, meaning flat, not dropped.

\paragraph{Two parameters that do not do what their labels say.}
Two details are worth recording because they are invisible in the formulas and
visible in the results. In CTA 3.0, the volatility used to \emph{size} the
position in \eqref{eq:cta3} has centre of mass $s$, the slow crossover length
--- not $\nu$, the value the ``Volatility'' slider sets. That slider reaches
only the price filter \eqref{eq:voladj}. So in CTA 3.0 the slow parameter does
double duty, and the shape of its parameter surface is not comparable with
that of its neighbours for that reason. In CTA 5.0 the crossover lengths are
fixed at $(32, 96)$ in the driver and are not search dimensions at all; what
is searched there is $(\nu, \kappa, \lambda)$.

\paragraph{Numerical conventions.}
The panel is read from a CSV of daily prices with hashed column names,
interpolated across internal gaps, and cast to \code{Float64}. Returns are
simple returns for the volatility estimates in \eqref{eq:cta2},
\eqref{eq:cta3}, \eqref{eq:cta4} and \eqref{eq:cta5}, and log returns for the
price filter \eqref{eq:voladj}. Covariance and correlation are the plain
exponentially weighted sample estimators; the only regularisation anywhere is
the shrinkage \eqref{eq:shrink}. Signals are computed with
Polars \citep{polars} expressions and the per-day solve of \eqref{eq:cta5}
with NumPy \citep{harris2020}.

\section{Results}
\label{sec:results}

All numbers come from the panel shipped with the repository:
$13{,}909$ daily rows from 1970-01-02 to 2023-04-26, $49$ assets with hashed
identifiers, between $3{,}556$ and $13{,}908$ observations each (median
$9{,}467$). The annualisation factor implied by the date index is
$260.70$ periods per year. Sharpe ratios are gross of costs and computed on
the portfolio's daily returns unless stated otherwise; ``turnover'' is annual
one-way turnover as a multiple of AUM, that is, the sum of absolute position
changes over all assets and days divided by AUM and by the number of years.

\subsection{The headline sequence}

\begin{table}[t]
\centering
\caption{The five systems at their notebook default parameters. Sharpe ratios
are gross of costs. ``Params'' counts the sliders the notebook exposes.}
\label{tab:headline}
\begin{tabular}{lrrrrrrr}
\toprule
System & Sharpe & Ann.\ vol. & Max DD & Kurtosis & Skew & Turnover & Params \\
\midrule
CTA 1.0 & $0.5606$ & $14.41\%$ & $-61.44\%$ & $30.54$ & $+0.49$ & $4.2$   & 2 \\
CTA 2.0 & $0.8794$ & $18.10\%$ & $-55.47\%$ & $6.44$  & $-0.41$ & $86.8$  & 3 \\
CTA 3.0 & $0.8776$ & $13.89\%$ & $-45.55\%$ & $7.99$  & $-0.41$ & $27.7$  & 4 \\
CTA 4.0 & $1.0712$ & $15.51\%$ & $-38.82\%$ & $3.94$  & $-0.36$ & $64.3$  & 4 \\
CTA 5.0 & $1.4652$ & $18.58\%$ & $-35.30\%$ & $1.44$  & $-0.17$ & $147.5$ & 6 \\
\bottomrule
\end{tabular}
\end{table}

Table~\ref{tab:headline} is the sequence as the notebooks present it, plus the
columns they do not print. Three readings of it are worth separating.

The Sharpe ratio rises by a factor of $2.6$ end to end, but in four steps of
which one is not a step: $+0.319$, $-0.002$, $+0.194$, $+0.394$. CTA 3.0 is
flat against CTA 2.0 --- $0.8776$ against $0.8794$ --- and marginally worse.
Anyone reading the printed Sharpe ratio at the bottom of each notebook sees
the third refinement do nothing.

The risk columns tell a cleaner story than the return column, and they are
monotone where the Sharpe ratio is not. Maximum drawdown improves at every
single step, from $-61\%$ to $-35\%$. Excess kurtosis falls from $30.5$ to
$1.4$: CTA 1.0's daily return distribution is dominated by rare enormous days
--- the inevitable consequence of a $\pm 1$ signal on $49$ assets sized
without reference to their volatility --- and by CTA 5.0 the distribution is
close to Gaussian. Skewness moves from $+0.49$ to $-0.17$, which is the
familiar trade: the crude system is a lottery ticket, the refined one is an
insurance seller. Both facts matter for how the Sharpe ratio should be read at
all, since its estimation error grows with both \citep{lo2002}: the same
$0.56$ measured on a kurtosis-$30$ series and on a kurtosis-$1.4$ series are
not equally trustworthy numbers.

And turnover rises by a factor of $35$. This is the column that reorders the
table, and it is the subject of the next section.

\subsection{Costs reorder the sequence}

Turnover in Table~\ref{tab:headline} is not monotone, and its pattern is the
inverse of what the notebooks suggest. CTA 1.0's own commentary calls the sign
function ``very expensive to trade as position changes are too extreme'',
which is true of each individual trade and false of the total: the sign
function is piecewise constant, so it trades only on a crossover flip, and its
$4.2$ turnovers a year make it by a wide margin the \emph{cheapest} of the
five. Volatility scaling is what makes trading continuous --- the denominator
of \eqref{eq:cta2} moves daily --- and it multiplies turnover by twenty, to
$86.8$.

\begin{table}[t]
\centering
\caption{Sharpe ratio net of a linear one-way trading cost, charged as
$\text{cost}_t = \text{turnover}_t \times \text{bps}/10^4$ on the daily
portfolio return.}
\label{tab:cost}
\begin{tabular}{lrrrrrr}
\toprule
System & $0$ bps & $1$ bps & $2$ bps & $5$ bps & $10$ bps & $20$ bps \\
\midrule
CTA 1.0 & $0.5606$ & $0.5576$ & $0.5547$ & $0.5459$ & $0.5312$ & $0.5018$ \\
CTA 2.0 & $0.8794$ & $0.8313$ & $0.7832$ & $0.6387$ & $0.3981$ & $-0.0776$ \\
CTA 3.0 & $0.8776$ & $0.8577$ & $0.8378$ & $0.7779$ & $0.6780$ & $0.4780$ \\
CTA 4.0 & $1.0712$ & $1.0297$ & $0.9881$ & $0.8632$ & $0.6551$ & $0.2404$ \\
CTA 5.0 & $1.4652$ & $1.3849$ & $1.3045$ & $1.0633$ & $0.6631$ & $-0.1174$ \\
\bottomrule
\end{tabular}
\end{table}

Table~\ref{tab:cost} charges a linear cost against turnover and recomputes the
Sharpe ratio. The step that looked like nothing is the only step in the
sequence that \emph{reduces} turnover, from $86.8$ to $27.7$, and the reason
is in \eqref{eq:cta3}: bounding the signal with $\tanh$ and standardising the
price with \eqref{eq:voladj} removes exactly the large signal excursions that
generate large position changes. Net of $5$ basis points, CTA 3.0 beats
CTA 2.0 by $+0.139$ Sharpe; net of $10$, by $+0.280$. The break-even is at
$0.06$ basis points --- effectively, CTA 3.0 dominates CTA 2.0 at any cost
above zero. The refinement is worth what it is worth entirely in trading
costs, and the headline number cannot see it.

\begin{table}[t]
\centering
\caption{One-way cost, in basis points of traded notional, at which the net
Sharpe ratios of the row and column systems coincide. ``---'' on an
off-diagonal entry means no crossing below $40$ bps: CTA 4.0 dominates
CTA 2.0 at every cost level.}
\label{tab:breakeven}
\begin{tabular}{lrrrrr}
\toprule
& CTA 1.0 & CTA 2.0 & CTA 3.0 & CTA 4.0 & CTA 5.0 \\
\midrule
CTA 1.0 & --      & $7.05$   & $18.60$ & $13.20$ & $11.72$ \\
CTA 2.0 & $7.05$  & --       & $0.06$  & ---     & $18.64$ \\
CTA 3.0 & $18.60$ & $0.06$   & --      & $8.94$  & $9.75$  \\
CTA 4.0 & $13.20$ & ---      & $8.94$  & --      & $10.21$ \\
CTA 5.0 & $11.72$ & $18.64$  & $9.75$  & $10.21$ & --      \\
\bottomrule
\end{tabular}
\end{table}

Table~\ref{tab:breakeven} generalises this into the full ordering. Two entries
deserve attention. Above $7.05$ basis points, the ten-line original beats
volatility scaling; above $11.72$, it beats the correlation-optimised finale,
and above $18.60$ it beats every other system in the sequence. And beyond
roughly $9$ to $10$ basis points, CTA 3.0 --- the step that appeared to
contribute nothing --- beats both of the optimised systems that follow it.

How much this matters depends on the cost of the panel, which the anonymised
data does not let us estimate. For liquid futures a one-way cost of one to
three basis points of notional is a defensible order of magnitude, and in that
range the gross ranking survives: at $2$ bps the five systems are still
ordered $0.55 < 0.78 \approx 0.84 < 0.99 < 1.30$. The honest summary is not
that the sequence is wrong but that its margin is thinner than the gross
numbers suggest, that it is monotone in cost-adjusted terms only below about
$9$ bps, and that a strategy trading $147$ times its AUM a year is a different
kind of object from one trading $4$ times, whatever their Sharpe ratios.

\subsection{Structure against parameters}

The notebooks state a thesis about parameter tuning: that ``such fundamental
flaws are not addressed by parameter-hacking or pimp-my-trading-system
steps''. This is testable, and it holds --- with a margin large enough to be
worth quantifying.

\begin{table}[t]
\centering
\caption{Sharpe ratio over the full grid of $852$ admissible crossover pairs
$(f, s)$, $f \in \{4, 8, \ldots, 96\}$, $s \in \{f+4, \ldots, 192\}$, with all
other parameters at their defaults. The last column is the percentile of the
notebook default $(32, 96)$ within that grid.}
\label{tab:grid}
\begin{tabular}{lrrrrrrrr}
\toprule
System & Min & $q_{25}$ & Median & $q_{75}$ & Max & SD & $\argmax$ & Default pct. \\
\midrule
CTA 1.0 & $0.3035$ & $0.4330$ & $0.4892$ & $0.5648$ & $0.6433$ & $0.0774$ & $(4, 64)$ & $72.5\%$ \\
CTA 2.0 & $0.6030$ & $0.7453$ & $0.8183$ & $0.8816$ & $1.0016$ & $0.0857$ & $(4, 84)$ & $73.6\%$ \\
CTA 3.0 & $0.6356$ & $0.7706$ & $0.8333$ & $0.8859$ & $0.9402$ & $0.0717$ & $(4, 76)$ & $70.7\%$ \\
CTA 4.0 & $0.7759$ & $0.9183$ & $1.0104$ & $1.0927$ & $1.2260$ & $0.1059$ & $(4, 76)$ & $67.4\%$ \\
\bottomrule
\end{tabular}
\end{table}

Table~\ref{tab:grid} sweeps the two parameters every practitioner tunes over
their entire admissible grid. CTA 5.0 is absent because its crossover lengths
are fixed in the driver and are not a search dimension there.

The central comparison is between the maximum of one row and the default of a
later one. CTA 1.0's best pair over all $852$ --- $0.6433$ at $(4,64)$ --- is
below CTA 2.0's median ($0.8183$), let alone its default ($0.8794$). CTA 2.0's
best ($1.0016$) is below CTA 4.0's median ($1.0104$) and default ($1.0712$).
CTA 4.0's best ($1.2260$) is below CTA 5.0's default ($1.4652$). Exhaustively
tuning the two parameters of any system in this sequence buys less than
adopting the structure of the system two steps later, and in the first case
buys less than adopting the structure one step later.

Two further readings. First, the hand-set defaults are good but not tuned:
$(32, 96)$ sits between the $67$th and $74$th percentile of each grid, which
is roughly where a sensible prior lands and well short of the maximum ---
consistent with sliders set by judgement rather than fitted. Second, every
grid is maximised at the fastest admissible $f = 4$, which on daily data is a
four-day centre of mass. That is a warning rather than a recommendation: the
fastest crossover is also the most expensive to trade, and
Table~\ref{tab:grid} is gross of costs.

\begin{table}[t]
\centering
\caption{Spearman rank correlation of the Sharpe ratio between systems, over
the $852$ shared grid points.}
\label{tab:rank}
\begin{tabular}{lrrr}
\toprule
& CTA 2.0 & CTA 3.0 & CTA 4.0 \\
\midrule
CTA 1.0 & $+0.978$ & $+0.961$ & $+0.966$ \\
CTA 2.0 &          & $+0.925$ & $+0.951$ \\
CTA 3.0 &          &          & $+0.963$ \\
\bottomrule
\end{tabular}
\end{table}

Table~\ref{tab:rank} explains why. The parameter surfaces of the four systems
are almost the same surface: the rank correlation between any pair is between
$0.925$ and $0.978$. Whatever the structural refinements are doing, they are
not changing which crossover lengths work. They shift the level of the whole
surface and leave its shape essentially intact. This is the precise sense in
which structure and parameters are not substitutes: a parameter search moves
you along a ridge, and each structural step lifts the entire ridge.

\begin{table}[t]
\centering
\caption{Optuna TPE search (seed $42$) over each system's own space, with the
repository's default trial budgets. The final column is the structural gain
from the preceding system at default parameters, for comparison.}
\label{tab:optuna}
\begin{tabular}{lrrrlr}
\toprule
System & Trials & Baseline & Best & Best parameters & Structural step \\
\midrule
CTA 1.0 & $200$ & $0.5606$ & $0.6208$ & $f{=}20$, $s{=}52$ & --- \\
CTA 2.0 & $200$ & $0.8794$ & $1.0263$ & $f{=}4$, $s{=}84$, $\nu{=}16$ & $+0.319$ \\
CTA 3.0 & $150$ & $0.8776$ & $1.0000$ & $f{=}48$, $s{=}56$, $\nu{=}4$ & $-0.002$ \\
CTA 4.0 & $150$ & $1.0712$ & $1.3042$ & $f{=}4$, $s{=}80$, $\nu{=}12$ & $+0.194$ \\
CTA 5.0 & $40$  & $1.4652$ & $1.5936$ & $\nu{=}12$, $\kappa{=}280$, $\lambda{=}0.8$ & $+0.394$ \\
\bottomrule
\end{tabular}
\end{table}

Table~\ref{tab:optuna} repeats the exercise with the repository's own Optuna
driver, searching the full space of each system rather than two of its
parameters. The searches find $+0.06$ to $+0.23$ Sharpe. The three structural
steps that pay find $+0.32$, $+0.19$ and $+0.39$. Only for CTA 4.0 does the
search exceed the structural step that produced the system it is searching,
and even there its best ($1.3042$) remains below CTA 5.0's untuned default.

\subsection{How much of the search gain is selection?}
\label{sec:selection}

A best-of-$K$ Sharpe ratio is biased upwards, and the size of the bias is the
question \citet{bailey2014} and \citet{harvey2015} exist to answer. The
standard correction subtracts the expected maximum of $K$ draws,
\begin{equation}
  \mathbb{E}[\max\nolimits_K] \approx \sigma \left[
    (1-\gamma_e)\Phi^{-1}\!\left(1 - \tfrac{1}{K}\right)
    + \gamma_e \Phi^{-1}\!\left(1 - \tfrac{1}{Ke}\right) \right],
\end{equation}
with $\gamma_e$ the Euler--Mascheroni constant, which for $K = 852$ gives a
factor of $3.209$ on $\sigma$. Which $\sigma$ is the whole difficulty.

Taking $\sigma$ as the cross-sectional standard deviation of the grid Sharpe
ratios, as \citet{bailey2014} prescribe, gives haircuts of $0.248$, $0.275$,
$0.230$ and $0.340$ for the four systems. Every one of these exceeds the
entire observed distance from the grid median to the grid maximum --- $0.154$
for CTA 1.0, $0.183$ for CTA 2.0. Read literally, the correction says the
best point of each grid is \emph{below} what pure selection would produce from
identical strategies, which is not a coherent conclusion. It is the expected
one: the correction assumes $K$ independent trials, and
Table~\ref{tab:rank} has just shown that neighbouring grid points are near
duplicates of each other. The effective number of independent trials in a
$852$-point grid over two smoothing lengths is far smaller than $852$, and
nothing here estimates it.

Taking $\sigma$ instead as the sampling standard deviation of a single Sharpe
ratio, $\sqrt{(1 + \mathrm{SR}^2/2)/T}$ with $T = 13{,}909$
\citep{lo2002}, gives haircuts of $0.029$ to $0.034$. That is a lower bound in
the opposite direction, since it ignores the parameter search entirely.

The useful statement is the bracket. The selection bias in
Table~\ref{tab:optuna} is somewhere between $0.03$ and $0.34$ Sharpe, and the
structural steps of $+0.32$, $+0.19$ and $+0.39$ are large relative to the
lower end and comparable to the upper end. This does not overturn
Section~\ref{sec:results}'s conclusion --- the structural steps are measured
at fixed default parameters, so they carry no selection bias of their own ---
but it does mean the search gains in Table~\ref{tab:optuna} should be read as
upper bounds on what a search would deliver out of sample, and quite possibly
as nothing at all.

\section{Limitations}

Everything above is in-sample in the strongest sense: one panel, one
realisation, all statistics computed on the same data that produced the
signals, and no walk-forward anywhere. The Sharpe ratios in
Table~\ref{tab:headline} are properties of a construction applied to a
history, not estimates of what any of these systems would earn. The panel runs
from 1970, and a trend-following strategy evaluated over that window is
evaluated over the period in which trend following is known to have worked
\citep{hurst2017}; CTA 1.0's own commentary observes that the system ``lost
its mojo in 2009'', and none of the statistics here are computed
sub-period.

The cost model of Table~\ref{tab:cost} is linear in traded notional with a
single rate across $49$ instruments, which is a crude proxy for a market
impact that is neither linear nor uniform. It is used to rank, not to price.
The break-even costs in Table~\ref{tab:breakeven} inherit that crudeness.

The assets are anonymised, so nothing here can be conditioned on asset class,
liquidity or contract size, and the cross-sectional constructions of
\eqref{eq:cta4} and \eqref{eq:cta5} cannot be checked for the concentrations
they might be taking. Asset counts vary over the panel --- the earliest rows
carry a fraction of the $49$ --- and \eqref{eq:cta5} handles that by
restricting to live assets each day, so its effective universe grows through
the sample.

No expected returns are estimated anywhere, no constraints on individual
positions or groups are imposed, and no costs enter any objective. Adding a
cost term to \eqref{eq:cta5}, which is where Table~\ref{tab:cost} suggests the
next improvement lies, is exactly the regularisation the last notebook lists
as future work and does not implement.

\section{Conclusion}

The ten-line CTA is a good pedagogical object because each of its flaws maps
to one identifiable repair, and the repairs are separable enough to be
measured one at a time. Measured that way, on this panel: the sequence's Sharpe
ratio gain is real and large, and it is not evenly distributed across the four
steps; the risk properties improve monotonically where the Sharpe ratio does
not, with excess kurtosis falling from $30.5$ to $1.4$; and the third step, the
one whose Sharpe ratio contribution is zero, is the one that makes the
sequence affordable, cutting turnover by two thirds and dominating the step it
replaces at any positive trading cost.

The claim the notebooks make about parameter tuning holds and holds strongly.
The best of $852$ crossover pairs never reaches the untuned default of the
system two steps later, and the parameter surfaces of the five systems are
rank-correlated above $0.92$ --- the refinements lift the surface rather than
reshaping it. Set against the selection bias that any best-of-$K$ number
carries, the case for spending effort on structure rather than search is about
as clean as a single panel can make it.

What the sequence does not show, and what Table~\ref{tab:breakeven} makes
uncomfortable, is that more machinery means more trading. CTA 5.0 turns over
$147$ times its AUM a year against CTA 1.0's $4$, and at a one-way cost above
about $12$ basis points the ten lines win. The next refinement in this
sequence should not be a better estimate of the correlation matrix. It should
be a cost term in \eqref{eq:cta5}.

\appendix

\section{Reproducibility}
\label{app:repro}

Every number above was produced from the repository at the commit that carries
this paper, with the price panel in
\code{book/marimo/notebooks/public/Prices\_hashed.csv} and the project's
locked environment. The five portfolios are built by the driver's own
builders, so the strategy code is the notebooks':

\begin{verbatim}
import sys; sys.path.insert(0, "book/marimo/notebooks")
import numpy as np, optimize as O

BUILD = {1: lambda: O.build_exp1(fast=32, slow=96),
         2: lambda: O.build_exp2(fast=32, slow=96, volatility=32),
         3: lambda: O.build_exp3(fast=32, slow=96, vola=32, clip=4.2),
         4: lambda: O.build_exp4(fast=32, slow=96, vola=32, clip=4.2),
         5: lambda: O.build_exp5(vola=32, clip=4.2, corr=200, shrinkage=0.5)}

for k, build in BUILD.items():
    p = build(); s = p.stats
    ppy = s.periods_per_year
    to = p.turnover["turnover"].to_numpy()
    print(k, float(s.sharpe()["returns"]), float(s.volatility()["returns"]),
          float(s.max_drawdown()["returns"]), float(s.kurtosis()["returns"]),
          float(s.skew()["returns"]), np.nansum(to) / (len(to) / ppy))
\end{verbatim}

Table~\ref{tab:cost} charges the cost through the same accounting, and
Table~\ref{tab:breakeven} bisects the difference of two such curves on a
$0.01$ bps grid over $[0, 40]$:

\begin{verbatim}
r = p.returns["returns"].to_numpy(); to = p.turnover["turnover"].to_numpy()
def net_sharpe(bps):
    x = r - to * bps / 1e4
    return x.mean() / x.std(ddof=1) * np.sqrt(ppy)
\end{verbatim}

Table~\ref{tab:grid} sweeps the grid

\begin{verbatim}
[(f, s) for f in range(4, 100, 4) for s in range(f + 4, 196, 4)]
\end{verbatim}

through the same builders with the remaining parameters at their defaults, and
Table~\ref{tab:rank} is \code{scipy.stats.spearmanr} on the resulting value
vectors. Table~\ref{tab:optuna} is the repository's driver verbatim:

\begin{verbatim}
python book/marimo/notebooks/optimize.py --experiment all --seed 42
\end{verbatim}

The searches are seeded, so that table reproduces exactly; the grid and the
cost curves are deterministic. The pinned Sharpe ratios in
\code{tests/expected\_sharpe.py} agree with column two of
Table~\ref{tab:headline} to all printed digits, which is the check that the
notebooks and the driver have not drifted apart.

\bibliographystyle{plainnat}
\bibliography{references}

\end{document}
